% ============================================================================
% Fichier  : front/01_avant_propos.tex (Avant-propos)
% Version  : 1.0 -- novembre 2024 (proposition à relire par l'auteur)
% ============================================================================

\chapter*{Avant-propos}
\addcontentsline{toc}{chapter}{Avant-propos}
\markboth{Avant-propos}{Avant-propos}

Une microfinance de Bafoussam perd quarante-sept millions de francs CFA en
une nuit. Trois mois plus tôt, l'éditeur de sa passerelle d'accès à distance
avait publié un correctif pour la faille que les attaquants ont empruntée.
L'enquête qui suit doit dire ce qui s'est passé, qui l'a fait si on peut le
savoir, et par où ils sont entrés. La dernière question renvoie à une autre
discipline, l'analyse de vulnérabilités, qui aurait pu éviter la première.

Ce cours réunit ces deux disciplines parce que, sur le terrain, elles se
répondent. L'investigation numérique part d'une trace et remonte vers une
cause~; la gestion des vulnérabilités part d'une faiblesse connue et cherche
à la fermer avant qu'elle ne laisse une trace. L'une sans l'autre laisse un
travail à moitié fait~: une enquête qui ne débouche sur aucune correction
prépare l'incident suivant, et un programme de correctifs qui ignore ce que
les enquêtes révèlent corrige au hasard.

Il s'adresse aux étudiants du Master Professionnel en Sécurité des Systèmes
Numériques de la Chaire UNESCO d'Afrique centrale, et aux praticiens qui
exercent ou exerceront comme analystes forensiques, analystes en réponse à
incident, analystes en gestion des vulnérabilités ou RSSI. Il suppose les
bases de la sécurité des systèmes. Il laisse au cours d'Ethical Hacking et
Pentesting (UE MSN-2S3-SEC03) l'exploitation des failles, et au cours de
droit du numérique (UE MSN-2S3-JUR01) la protection des données
personnelles.

Trois partis pris le traversent. Le premier est de tenir la preuve pour
fragile~: une copie mal faite, un horodatage mal lu ou un scellé mal tenu
suffisent à faire écarter une constatation exacte. Le deuxième est de
séparer, à chaque page, ce que l'on a constaté de ce que l'on en déduit, et
de dire avec quelle confiance on le déduit. Le troisième est de coter ce que
l'on ne sait pas~: l'instrument \OUI{}~/ \NON{}~/ \NEK{} du programme de
recherche de l'auditologie sert ici à dire, pour chaque conclusion
d'enquête, si elle est établie, réfutée, ou hors de portée des preuves
disponibles.

Les deux organisations qui accompagnent le lecteur, NYANGA CRÉDIT et
SOLITRANS LOGISTIQUE, sont fictives, comme leurs dirigeants, leurs
prestataires, les dates, les adresses et les montants. La faille de la
passerelle porte un identifiant fictif. Les mécanismes, les outils, les
normes et les textes cités sont exacts. Les articles de loi sont cités dans
leur texte officiel~; les normes sous licence sont paraphrasées.

Le livre est accompagné d'un livret des corrigés, qui explique et annote
chaque exercice, et de NDAP SUKULU, où les fiches de questions et les dix
laboratoires du cours permettent de pratiquer chaque chapitre.

\bigskip
\hfill\begin{minipage}{0.5\textwidth}
\raggedleft
\auteurcours\\
\laboratoire, Yaoundé, novembre 2024
\end{minipage}
